\section{Lesson 6: Demonstration and Its Limits}

Now that you can pull apart a metaphysical paragraph---naming its claim, fixing its terms, sorting the warrant of each premise, isolating its burden, bounding its consequence, and classifying what denial costs---you are finally ready to test the highest claim a classical philosophical text can make for itself: that an argument is a \textbf{demonstration}.

To a modern ear, ``demonstration'' usually means nothing more than a logically valid argument whose premises are true. In Aristotle's \emph{Posterior Analytics} and St.\ Thomas Aquinas's commentary upon it, validity is the cheapest part of a demonstration. An argument can follow every rule of syllogistic form, every single premise inside it can be true, and it can still fail to be a demonstration. Everything demanding in classical proof lies in the status of the premises and in the exact causal work done by the middle term.

Before stepping into those historical standards, keep the boundary between our reading method and the ancient theory completely clear. The six questions we have been practising---Claim, Meaning, Warrant, Burden, Consequence, and Denial---along with our eight kinds of warrant and five kinds of denial, are a modern reader's diagnostic tool, never an apparatus Aristotle or Aquinas themselves used. Where our eight kinds of warrant sort whatever an author happens to put on the page (including revelation and bare provenance), Aristotle's conditions on demonstrative premises (\emph{Posterior Analytics} I.2, 71b17--72a8) define what a completed science (\emph{episteme}, Latin \emph{scientia}) must achieve: its premises must be \emph{true, primary, immediate, better known than the conclusion, prior to it, and causes of it}. A second crucial difference in vocabulary separates the two: what we have called a ``metaphysical principle'' is any substantive starting claim about reality, contested or not; for Aristotle, a \emph{principle} (\emph{arche}, Latin \emph{principium}, \emph{Post.\ An.}\ I.2, 72a8--25) is strictly an immediate, necessary proposition with nothing prior to it. Honoring those differences lets us learn from Aristotle and Aquinas without putting our modern questions into their mouths.

\subsection{The Middle Term as Cause: \emph{Quia} and \emph{Propter Quid}}

Scientific knowledge (\emph{episteme}), Aristotle writes in \emph{Posterior Analytics} I.2 (71b8--16; cf.\ \emph{Nicomachean Ethics} VI.3, 1139b18--22), is knowing the cause on which a fact depends, knowing it as the cause of that fact and of no other, and knowing that the fact cannot be otherwise. A demonstration (\emph{apodeixis}) is a syllogism productive of that kind of knowledge.

That requirement transforms what the \textbf{middle term} of a syllogism must be. In formal logic, a middle term is simply whatever shared term links minor and major premises validly. In a strict demonstration of the reasoned fact---\emph{demonstratio propter quid} (``demonstration on account of which'')---the middle term must be the \textbf{cause} in virtue of which the predicate belongs to the subject (\emph{Post.\ An.}\ I.2, 71b17--72a8; II.11, 94a20--95a9). Aristotle permits any of the four causes---definable form, necessitating antecedent, efficient cause, or final cause---to serve as middle term. In his literal commentary (\emph{Expositio libri Posteriorum Analyticorum} II, lect.\ 1, n.\,8; cf.\ lect.\ 9, nn.\,2--3), Aquinas condenses the rule to one line: ``the cause is the middle in a demonstration, which makes us know; for to know is to know the cause of a thing.''

Watch what happens when one conclusion is proved in opposite directions, as in Aristotle's planets example from Lesson 0 (\emph{Post.\ An.}\ I.13, 78a22--79a16) and Aquinas's \emph{duplex est demonstratio} (\emph{ST} I, q.\,2, a.\,2, corp.). In the \textbf{\emph{quia} syllogism}---whatever does not twinkle is near; the planets do not twinkle; therefore the planets are near---the middle term is \emph{not twinkling}, and the inference runs \emph{from effect to cause}, from what is better known to our senses toward what is prior in nature. Because nearness and steady light convert reciprocally among naked-eye celestial bodies, the syllogism genuinely demonstrates \emph{that} the planets are near without explaining \emph{why}. Reverse the middle and major terms to form the \textbf{\emph{propter quid} syllogism}---whatever is near does not twinkle; the planets are near; therefore the planets do not twinkle---and the middle term becomes \emph{being near}, the physical cause, running \emph{from cause to effect} to show both \emph{that} and \emph{why}.

Guard against one distortion: \emph{demonstratio quia} is not a failed proof or an inferior guess. For both Aristotle and Aquinas, \emph{quia} is a genuine mode of demonstration whenever an effect is better known to us than its cause. What the distinction prevents is sliding silently from proving \emph{that} something must exist to claiming you have explained \emph{why} it is what it is.

\subsection{Demonstration Is Not Explanation}

A second distinction sits directly beside \emph{quia} and \emph{propter quid}: \textbf{demonstration is not the same thing as explanation.} An explanation renders a puzzling phenomenon intelligible; a demonstration establishes a conclusion from premises meeting strict logical and causal conditions. Because \emph{demonstratio propter quid} does both at once, readers constantly blur them---yet they diverge in both directions. A \emph{demonstratio quia} demonstrates without explaining (smoke proves \emph{that} a fire burns behind a ridge without explaining what ignited it). Conversely, an illuminating hypothesis can unify scattered facts and dissolve a puzzle without ever demonstrating that its proposed cause exists or could not be otherwise. Inference to the best explanation belongs under \textbf{hypothesis}, not causal demonstration: \textbf{feeling illuminated is not proof.}

\subsection{How a Valid, Sound Argument Fails to Be a Demonstration}

Even when an argument is valid and all its premises are true, six specific defects keep it from being a strict Aristotelian demonstration. You do not need to memorize a count; recognize each when a page presents it:

\begin{itemize}
\item \textbf{Premises less known in nature.} Reasoning from posterior effects or symptoms to a prior cause (\emph{Post.\ An.}\ I.2, 71b29--72a5; I.13, 78a22--38) yields at best \emph{demonstratio quia}, since strict \emph{propter quid} demonstration starts from what is better known in nature.
\item \textbf{A non-causal middle term.} A true accidental companion or overly remote feature neither causes nor flows properly from the predicate (\emph{Post.\ An.}\ I.4, 73a34--b26; I.13, 78b13--34: arguing that a wall does not breathe because it is not an animal uses a middle too remote, since not every animal breathes).
\item \textbf{Merely probable premises.} Reasoning from \emph{endoxa}---reputable opinions held by everyone, by most, or by the wise (\emph{Topics} I.1, 100a18--b23)---is \emph{dialectical} argument, not demonstration.
\item \textbf{Revealed premises.} In \emph{ST} I, q.\,1, a.\,2, corp., Aquinas explains that sacred doctrine is a science proceeding from principles known in the light of a higher science---the science of God and the blessed---as perspective proceeds from geometry and music from arithmetic. Truths such as ``God exists'' that natural reason can also reach are \emph{praeambula fidei} (``preambles to the articles,'' \emph{ST} I, q.\,2, a.\,2, ad 1; cf.\ \emph{SCG} I, c.\,3), which one who cannot follow a philosophical proof may rightly hold by faith. Yet \textbf{a proposition's being both naturally reachable and revealed in Scripture never turns a theological argument from Scripture into a philosophical demonstration}: appealing to Exodus 3:14 argues from revelation (canonical Scripture, \emph{ST} I, q.\,1, a.\,8, ad 2), not from natural first principles.
\item \textbf{A non-necessary conclusion.} Strict scientific knowledge requires necessary premises; there is no Aristotelian demonstration of chance conjunctions or passing accidents (\emph{Post.\ An.}\ I.6, 74b5--75a17; I.8, 75b21--36; I.30, 87b19--27).
\item \textbf{Crossing subject genera.} Demonstration stays inside its subject genus (\emph{Post.\ An.}\ I.7, 75a38--b20), barring subordinate sciences such as optics under geometry (I.13, 78b34--79a16). Thus Aquinas separates metaphysics (\emph{ens commune}) from physics and mathematics not by a neo-scholastic ``third degree of abstraction,'' but by \emph{separatio}---a negative judgment that being need not depend on matter (\emph{Super Boet.\ De Trin.}\ q.\,5, a.\,3).
\end{itemize}

Hold fast to one balancing judgment: \textbf{``not a demonstration'' never means ``worthless.''} Because demonstration can neither run in a circle nor regress forever, first principles themselves cannot be demonstrated (\emph{Post.\ An.}\ I.3, 72b5--73a20); Aristotle accounts for our grasp of them through \emph{epagoge} (induction from perception, memory, and experience) and \emph{nous} (\emph{Post.\ An.}\ II.19, 99b18--100b17; \emph{Nic.\ Eth.}\ VI.6, 1140b31--1141a8)---a chapter scholars still actively dispute, dividing over whether it describes a natural cognitive process or an epistemic capacity. Equally, \textbf{dialectic} (\emph{Topics} I.2, 101a25--37) is the indispensable path by which first principles are examined at all. Probable argument, dialectic, and explanation do indispensable work; accurate reading simply refuses to mislabel them.

\subsection{Worked Passage 1: Valid, Sound, and Not a Demonstration}

Consider a crisp syllogism built on a traditional human proper attribute (\emph{risibile}, the natural capacity for laughter, noted by Aquinas in \emph{ST} I, q.\,3, a.\,4, corp.\ and \emph{De ente et essentia} c.\,4):

\begin{quote}
Every animal naturally capable of laughter is a rational animal. Every human being is an animal naturally capable of laughter. Therefore, every human being is a rational animal.
\end{quote}

\textbf{Claim and Meaning.} Because every animal endowed with the natural capacity to laugh possesses reason, and all humans possess that natural aptitude (\emph{proprium}, not merely chuckling right now), all humans are rational animals. \textbf{Warrant and Burden.} The minor rests on \textbf{experience}; the major pairs \textbf{experience} with a \textbf{metaphysical principle} linking comprehension of humour exclusively to rational intellect---the exact \textbf{burden} an ethologist could contest.

\textbf{Consequence and Verdict.} Granting both premises, the syllogism is valid and sound---yet it is \textbf{not a strict Aristotelian \emph{demonstratio propter quid}}. Human beings are not rational \emph{because} they can laugh; they can laugh \emph{because} they are rational. Placing a downstream proper effect in the middle slot fails the causal-middle and priority-in-nature conditions (\emph{Post.\ An.}\ I.2, 71b17--72a8; I.13, 78a22--38), achieving at best \emph{demonstratio quia}. Contrast a still weaker valid, sound syllogism using a purely \textbf{accidental true middle}: ``Every feathered creature has a heart; every swan is a feathered creature; therefore every swan has a heart.'' Feathers neither cause a heart nor flow properly from having one, failing \emph{per se} mediation altogether. \emph{Not yet established:} \emph{why} humans are rational or \emph{why} swans have hearts. \textbf{Denial.} Rejecting the major premise incurs no \textbf{contradiction}, only a \textbf{rival framework} or empirical dispute in comparative ethology.

\subsection{Worked Passage 2: Self-Evident in Itself, Not Self-Evident to Us}

Return briefly to \emph{Summa Theologiae} I, q.\,2, a.\,1, corp.\ (examined in Lesson 3; cf.\ \emph{SCG} I, cc.\,10--11, esp.\ c.\,11, n.\,1), where Aquinas distinguishes what is self-evident in itself (\emph{per se notum quoad se}) from what is self-evident to us (\emph{per se notum quoad nos}):

\begin{quote}
Therefore I say that this proposition, ``God exists,'' of itself is self-evident, for the predicate is the same as the subject, because God is His own existence as will be hereafter shown (Q.\,3, Art.\,4). Now because we do not know the essence of God, the proposition is not self-evident to us; but needs to be demonstrated by things that are more known to us, though less known in their nature---namely, by effects.
\end{quote}

\textbf{Claim and Meaning.} Because God's nature is identical with His existence, ``God exists'' is objectively self-evident in itself; yet because human beings lack direct intellectual grasp of God's essence, it is not self-evident to us and must be proved from effects. Guard sharply against relativism: saying a proposition is self-evident \emph{quoad se} yet not \emph{quoad nos} does \emph{not} make self-evidence ``whatever a reader happens to find obvious.'' Whether predicate belongs to essence is an objective fact; what varies relative to us is our cognitive access.

\textbf{Warrant and Burden.} The rule of self-evidence is a \textbf{definition}; ``God is His own existence'' is a forward \textbf{prior result} deferred to q.\,3, a.\,4 (\textbf{provenance only} here); our lack of direct vision of God's essence rests on epistemological \textbf{experience}. The double \textbf{burden} is that God's essence is identical with His existence and that no human concept of ``God'' directly exhibits God's essence.

\textbf{Consequence and Denial.} Blocking self-evidence \emph{quoad nos} blocks conceptual shortcuts (Reply 2) and fixes the exact mode of proof in the next article (\emph{ST} I, q.\,2, a.\,2, ad 2): because we lack God's essence (\emph{quod quid est}), we must proceed by \emph{demonstratio quia} from effects, taking what the name signifies (\emph{quid significet nomen}) in place of the essence as middle term. Note candidly that scholars still dispute whether a \emph{quia} proof using a nominal middle meets Aristotle's strictest \emph{Post.\ An.}\ I.2 conditions for \emph{episteme} or extends them. \emph{Not yet established:} that God exists, that God's essence is His existence, or that any proof from effects succeeds. Denying either half of the burden commits to a \textbf{rival framework} (rejecting essence-existence identity, or defending an ontological proof), never bare \textbf{contradiction}.

\subsection{Worked Passage 3: Principle or Conclusion? Returning to the First Way}

Return a third time to \emph{ST} I, q.\,2, a.\,3, corp., focusing on one decisive demonstrative question from \emph{Posterior Analytics} I.2 (71b17--72a25): does ``whatever is in motion is put in motion by another'' stand as an \textbf{immediate first principle} (with no middle term prior to it) or as a \textbf{derived conclusion}?

Look at the single word ``for'' (\emph{enim}) right after the claim: ``Now whatever is in motion is put in motion by another, \emph{for} nothing can be in motion except it is in potentiality\dots\ But nothing can be reduced from potentiality to actuality, except by something in a state of actuality.'' Read as an \textbf{immediate first principle}, the sentence runs directly into trouble under Aristotle's strict conditions, since it is neither a verbal tautology nor immediately obvious across inertial motion or quantum decay. Read faithfully to Aquinas's syntax as a \textbf{derived conclusion}, it follows as a \textbf{logical consequence} of a compressed sub-argument combining a \textbf{definition} of motion, non-contradiction, and a deeper \textbf{metaphysical principle} (nothing reduces potency to act except an actual being).

That reading shifts the \textbf{burden} one layer deeper---onto the actualization principle and the reality of potency---rescuing Aquinas from the charge of assuming a bare slogan while showing that the First Way is at best a compressed \emph{demonstratio quia} resting on a contested starting point rather than an undisputed axiom. \emph{Not yet established:} that the First Way satisfies \emph{demonstratio propter quid}, or that natural movers form one unique cosmic \emph{per se} series. Denying the deeper actualization principle commits to the microstructural or Humean \textbf{rival frameworks} of Lesson 5, not \textbf{contradiction}.

\subsection{Exercises}

\subsubsection*{Exercise 6.1}
Read the following short argument carefully.

\begin{quote}
Every afternoon on which a properly calibrated mercury barometer drops rapidly at sea level is an afternoon of falling local atmospheric pressure. This afternoon our calibrated mercury barometer at sea level is dropping rapidly. Therefore, this afternoon local atmospheric pressure is falling.
\end{quote}

This argument is logically valid, and both of its premises are true. Is it a strict Aristotelian demonstration showing \emph{why} the conclusion holds? If not, name the exact condition it fails, and say what kind of valid reasoning it still achieves.

\subsubsection*{Exercise 6.2}
Compare the following pair of arguments, adapted from Aristotle's reciprocal-predication example in \emph{Posterior Analytics} II.16 (98b4--10; note that this illustration stands in Book II, not I.13):

\begin{quote}
Argument A: Every deciduous plant is broad-leaved. Every vine is deciduous. Therefore, every vine is broad-leaved.
\end{quote}

Now read the companion argument immediately below it:

\begin{quote}
Argument B: Every broad-leaved plant sheds its leaves in autumn (on Aristotle's botanical premise that broad leaves allow sap to coagulate at the leaf-stalk). Every vine is broad-leaved. Therefore, every vine sheds its leaves in autumn.
\end{quote}

Identify which argument shows \emph{that} (\emph{quia}) and which shows \emph{why} (\emph{propter quid}) on Aristotle's botanical premise. Name the middle term and direction of inference in each, explain what turns a \emph{quia} proof into a \emph{propter quid} proof---and why that reversal is unavailable when arguing from creatures to God.

\subsubsection*{Exercise 6.3}
Read the following valid, sound argument whose final clause makes a false demonstrative claim, alongside a second defective passage blending revelation and hypothesis:

\begin{quote}
Passage A: Every celestial body that turns dark copper-red at the time of a full moon is undergoing a lunar eclipse. Tonight's moon is turning dark copper-red at the time of a full moon. Therefore, tonight's moon is undergoing a lunar eclipse---which demonstrates \emph{why} the moon is eclipsed tonight.
\end{quote}

Compare that with a common conflation in theological apologetics:

\begin{quote}
Passage B: That the universe was created out of nothing is strictly demonstrated by Genesis 1:1---``In the beginning God created heaven and earth''---reinforced by the fact that a transcendent Creator offers the simplest explanation of mathematical order in nature. Because Scripture is infallible and no rival theory unifies cosmic order so cleanly, we possess a strict philosophical demonstration of creation out of nothing.
\end{quote}

For Passage A, diagnose the failure as \emph{propter quid}, lower its claim candidly, and then repair its middle term into a genuine causal \emph{propter quid} argument without smuggling the conclusion into a definition (verifying your repair with the hostile-reader test). For Passage B, sort the warrant of each strand and state what genuine standing each strand retains once ``strict philosophical demonstration'' is removed.

\subsubsection*{Exercise 6.4}
Here is an unannotated excerpt from St.\ Thomas Aquinas, \emph{Summa Theologiae} I, q.\,2, a.\,2 (``Whether It Can Be Demonstrated That God Exists?'', translated by the Fathers of the English Dominican Province, 1920). First, the body of the article:

\begin{quote}
Demonstration can be made in two ways: One is through the cause, and is called \emph{a priori}, and this is to argue from what is prior absolutely. The other is through the effect, and is called a demonstration \emph{a posteriori}; this is to argue from what is prior relatively only to us. When an effect is better known to us than its cause, from the effect we proceed to the knowledge of the cause. And from every effect the existence of its proper cause can be demonstrated, so long as its effects are better known to us; because since every effect depends upon its cause, if the effect exists, the cause must pre-exist. Hence the existence of God, in so far as it is not self-evident to us, can be demonstrated from those of His effects which are known to us.
\end{quote}

Immediately afterward, Aquinas answers Objection 2 (which objected that the essence is the middle term of demonstration, and we do not know God's essence):

\begin{quote}
When the existence of a cause is demonstrated from an effect, this effect takes the place of the definition of the cause in proof of the cause's existence. This is especially the case in regard to God, because, in order to prove the existence of anything, it is necessary to accept as a middle term the meaning of the word, and not its essence, for the question of its essence follows on the question of its existence. Now the names given to God are derived from His effects; consequently, in demonstrating the existence of God from His effects, we may take for the middle term the meaning of the word ``God''.
\end{quote}

Read both paragraphs with a pencil. Judge what kind of demonstration Aquinas is claiming for God's existence, identify what takes the place of the essence as the middle term and why, locate the burden-bearing premise on which this strategy turns, and state the limits Aquinas himself attaches to what such a demonstration can establish.

\subsection{Solutions}

\subsubsection*{Solution to Exercise 6.1}
No: although the syllogism is valid and both premises are true, it is \emph{not} a strict Aristotelian \emph{demonstratio propter quid}. The exact condition it fails is that its \textbf{middle term is not the cause} of the predicate's belonging to the subject (and its premises are posterior in nature to the conclusion). The mercury column falling in the glass tube is a mechanical \emph{effect} of reduced air weight pressing on the mercury reservoir, not the cause that makes atmospheric pressure drop. Because the instrument reading is a reliable effect better known to immediate sight, the syllogism achieves a genuine \textbf{\emph{demonstratio quia}}: it proves \emph{that} atmospheric pressure is falling without explaining \emph{why}.

\subsubsection*{Solution to Exercise 6.2}
\textbf{Argument A} shows only \emph{that} (\emph{demonstratio quia}): its middle term is \emph{being deciduous}, running from the observable effect (shedding leaves) back to the leaf structure without explaining why vines have broad leaves. \textbf{Argument B} shows \emph{why} (\emph{demonstratio propter quid}): on Aristotle's botanical premise, its middle term---\emph{being broad-leaved}---names the structural cause explaining why the vine sheds its leaves, reversing middle and major in reciprocal predication (\emph{Post.\ An.}\ II.16, 98b4--10). Turning \emph{quia} into \emph{propter quid} requires grasping the cause's nature closely enough to deduce the effect from it---a reversal unavailable in natural theology because we do not grasp the divine essence in this life and finite effects are not proportionate to an infinite cause (\emph{ST} I, q.\,2, aa.\,1--2, ad 3).

\subsubsection*{Solution to Exercise 6.3}
In \textbf{Passage A}, turning dark copper-red at full moon is an optical \emph{effect} (sunlight refracting through Earth's atmosphere into Earth's shadow), not the \emph{cause} of the eclipse, failing the causal-middle condition. Candidly lowered, the conclusion reads: ``Therefore, tonight's moon is undergoing a lunar eclipse, established \emph{quia} from its visible optical effect.'' Repaired into a genuine \emph{propter quid} proof, define ``lunar eclipse'' neutrally as \emph{a sudden darkening across the face of the full moon} and place the physical efficient cause in the middle slot: ``Whenever an opaque body lit solely by the sun has the Earth interposed directly between it and the sun, it undergoes a sudden darkening; tonight's full moon has the Earth so interposed; therefore tonight's full moon undergoes a sudden darkening (lunar eclipse).'' This repair passes the \textbf{hostile-reader test}: because ``lunar eclipse'' is defined strictly by the visible darkening rather than by ``having the Earth interposed,'' an opponent can accept every definition and still dispute the astronomical cause or tonight's alignment, proving no conclusion was smuggled into a definition.

In \textbf{Passage B}, two non-demonstrative strands are mislabelled. Genesis 1:1 rests strictly on \textbf{revelation}: even where revealed truth overlaps with \emph{praeambula fidei} (while remembering Aquinas holds in \emph{ST} I, q.\,46, a.\,2 that a \emph{temporal beginning} belongs to faith alone), citing Scripture is theological argument from revealed authority, not Aristotelian natural demonstration. Appealing to the simplest explanation of mathematical order rests on \textbf{hypothesis} (inference to the best explanation), which illuminates without proving necessity and can be contested by rival frameworks. Removing the overclaim belittles neither strand: Genesis 1:1 retains full authority within sacred doctrine, and cosmic order retains genuine explanatory weight.

\subsubsection*{Solution to Exercise 6.4}
Aquinas claims strictly \textbf{\emph{demonstratio quia}}---demonstration through effects prior relative to us (note that the 1920 Dominican translation renders \emph{propter quid} and \emph{quia} as \emph{a priori} and \emph{a posteriori}). In place of God's unknown essence (\emph{quod quid est}), Aquinas takes the effect---specifically \textbf{what the name ``God'' signifies} (\emph{quid significet nomen}), imposed from effects---as middle term, because \emph{whether} something exists (\emph{an est}) precedes \emph{what} it is (\emph{quid est}).

The \textbf{burden-bearing premise} is not the near-definitional line ``every effect depends on its cause'' (an opponent can grant that without taking a step toward God), but the substantive bridge \emph{that the sensible things better known to us are in fact effects of a transcendent proper cause answering to what ``God'' signifies}---whose denial is a \textbf{rival framework} (e.g., an uncaused physical cosmos), not contradiction. Finally, two limits stand \textbf{not yet established}: Reply 2 (with Reply 3) restricts the proof to \emph{that} God exists without grasping \emph{what} God is; and Article 2 argues only demonstrability in principle, supplying no actual proof yet (those begin in Article 3) and leaving open the scholarly dispute over whether a nominal-definition middle satisfies Aristotle's strictest \emph{Post.\ An.}\ I.2 conditions for \emph{episteme}.
